\textbf{Proof.}
For every \([C]\in\mathcal B_Z\), lem:zhao-rank-two-stratum-identification gives \(w^+(\iota(C))=w(C)\).  Since \(\iota(C)\) is feasible in the novel completion,
\[
 v_{\mathrm{mm}}^+\le\inf_{C\in\mathcal B_Z}w(C). \tag{1}
\]
Conversely, a rank-two completion class is uniquely \(\iota(C)\), while a singular completion class admits, by lem:rank-stratified-projector-reduction, ordinary classes \(C_j\) with \(\limsup_jw(C_j)\le w^+(B)\).  Hence \(\inf_{C\in\mathcal B_Z}w(C)\le w^+(B)\) for every completion class.  Taking the minimum over that class and combining with (1) proves
\[
 v_{\mathrm{mm}}^+=\inf_{C\in\mathcal B_Z}w(C).
\]
Prop:kappa-value-exhaustion gives the second equality.

If a rank-two completion minimizer equals \(\iota(C)\), prop:kappa-exhaustion supplies \(\kappa_C>0\) with \(C\in\mathcal B_{\kappa_C}\).  Then
\[
 v_{\mathrm{mm}}^+
 \le\min_{A\in\mathcal B_{\kappa_C}}w(A)
 \le w(C)=v_{\mathrm{mm}}^+,
\]
so \(C\in\mathcal A_{\mathrm{mm}}(\kappa_C)\) at a value-attaining threshold.  Conversely, every member of a fixed-threshold argmin whose value is \(v_{\mathrm{mm}}^+\) maps under \(\iota\) to a rank-two completion minimizer.  This proves the exact rank-two identity.  The final displayed set in the statement merely partitions the remaining equality set by rank; it retains possible singular ties and asserts no existence. \(\square\)